\documentclass[11pt]{article}
\usepackage[margin=1in]{geometry}
\usepackage{amsmath,amssymb,amsthm,mathtools}
\usepackage{graphicx}
\usepackage{xcolor}
\usepackage{microtype}
\usepackage[colorlinks=true,linkcolor=blue!55!black,citecolor=blue!55!black,urlcolor=blue!55!black]{hyperref}

\newtheorem{theorem}{Theorem}
\newtheorem{proposition}[theorem]{Proposition}
\newtheorem{lemma}[theorem]{Lemma}
\newcommand{\Cn}{\mathbb C^n}
\newcommand{\Fc}[2]{F^{#1}_{#2}}
\newcommand{\ip}[2]{\langle #1,#2\rangle}
\newcommand{\Sc}{\mathcal S}
\newcommand{\Tc}{\mathcal T}
\newcommand{\Tm}{T}

\title{Canonical-inner-product Toeplitz operators on Fock--Sobolev type spaces\\
\large An identity-plus-compact transfer theorem}
\author{Run \texttt{fa\_banach\_001}, lane 6}
\date{27 August 2026}

\begin{document}
\maketitle

\begin{center}
\fbox{\parbox{0.90\textwidth}{\textbf{Status.}
Candidate full result, likely valid.  Theorems A--C of
Cho--Isralowitz--Joo transfer to the canonical inner product, including the
canonical Berezin-transform characterization for every Schatten exponent
\(0<t<\infty\).  The proof is a direct operator-theoretic transfer, augmented
by an exact canonical-kernel localization lemma.}}
\end{center}

\section{Source question}

The source is H. R. Cho, J. Isralowitz and J.-C. Joo, \emph{Toeplitz
operators on Fock--Sobolev type spaces}, arXiv:1307.4829, later published in
Integral Equations and Operator Theory 82 (2015), 1--32 \cite{CIJ}.
Section 5 asks whether its Theorems A, B and C remain true when the Toeplitz
operator is defined from the reproducing kernel for the natural inner product,
and specifically asks for a direct operator-theoretic transfer between the two
equivalent Hilbert inner products.  The decisive source passage is reproduced
below.

\begin{figure}[h]
  \centering
  \includegraphics[width=0.96\textwidth]{figures/open_problem_crop.png}
  \caption{The open problem in Section 5, page 31 of arXiv:1307.4829v2.}
\end{figure}

\section{\mbox{\hspace{0.45em}Definitions and statement}}

Write
\[
 d v_\alpha(z)=(1+|z|)^{-\alpha}\,dv(z),\qquad
 \|f\|_{\Fc p\alpha}^p=\int_{\Cn}|f(z)|^p e^{-p|z|^2/2}\,dv_\alpha(z).
\]
The canonical Hilbert inner product on \(\Fc 2\alpha\) is
\[
 C(f,g):=\int_{\Cn}f\overline g\,e^{-|z|^2}(1+|z|)^{-\alpha}\,dv.
\]
Let \(K^c_z\) be its reproducing kernel.  For a positive locally finite
Borel measure \(\mu\), the canonical Toeplitz operator is initially defined on
polynomials by
\[
 \Tc_\mu f(z)=\int_{\Cn} f(w)\overline{K^c_z(w)}e^{-|w|^2}
 (1+|w|)^{-\alpha}\,d\mu(w).                                      \tag{2.1}
\]
Equivalently,
\[
 C(\Tc_\mu f,g)=C_\mu(f,g):=
 \int_{\Cn}f\overline g\,e^{-|z|^2}(1+|z|)^{-\alpha}\,d\mu.       \tag{2.2}
\]

Let \(M\) be the modified inner product used in \cite{CIJ}.  If
\(\alpha\leq0\),
\[
 M(f,g)=\int f\overline g\,e^{-|z|^2}|z|^{-\alpha}\,dv.
\]
If \(\alpha>0\), let \(P\) be the Taylor projection onto homogeneous
degrees at most \(\alpha/2\), let \(Q=I-P\), and put
\[
 M(f,g)=\int Pf\,\overline{Pg}\,e^{-|z|^2}\,dv
 +\int Qf\,\overline{Qg}\,e^{-|z|^2}|z|^{-\alpha}\,dv.           \tag{2.3}
\]
The source operator \(\Tm_\mu\) represents the analogous form \(M_\mu\).

Fix a sufficiently small radius \(r>0\) and an \(r\)-lattice
\(\{a_j\}\subset\Cn\).  The choice of lattice and radius is immaterial for
the conditions below.

\begin{theorem}[Canonical operator/geometric criteria]\label{thm:main}
Let \(\alpha\in\mathbb R\), and let \(\mu\geq0\) satisfy the natural
dense-domain condition for \((2.1)\).  Then:
\begin{enumerate}
\item For every \(1\leq p<\infty\), \(\Tc_\mu\) extends boundedly on
\(\Fc p\alpha\) if and only if
\[
 \sup_{z\in\Cn}\mu(B(z,r))<\infty.                                \tag{2.4}
\]
\item For every \(1\leq p<\infty\), \(\Tc_\mu\) is compact on
\(\Fc p\alpha\) if and only if
\[
 \mu(B(z,r))\longrightarrow0\quad (|z|\to\infty).                 \tag{2.5}
\]
\item For every \(0<t<\infty\), as an operator on the canonical Hilbert
space \(\Fc2\alpha\),
\[
 \Tc_\mu\in\Sc_t
 \quad\Longleftrightarrow\quad
 \mu(B(\mathord\cdot,r))\in L^t(dv)
 \quad\Longleftrightarrow\quad
 \{\mu(B(a_j,r))\}_j\in\ell^t.                                   \tag{2.6}
\]
\end{enumerate}
Moreover, if \(\widetilde\mu_c(z)=C(\Tc_\mu k^c_z,k^c_z)\) is the
canonical Berezin transform, then it is bounded exactly under \((2.4)\), and
it vanishes at infinity exactly under \((2.5)\).  For every \(0<t<\infty\),
\[
 \widetilde\mu_c\in L^t(dv)\quad\Longleftrightarrow\quad
 \Tc_\mu\in\Sc_t.                                                  \tag{2.7}
\]
\end{theorem}

Thus all clauses of source Theorems A, B and C survive the canonical inner
product.

\section{Proof intuition}

Equivalent Hilbert norms alone do not preserve Toeplitz operators, because a
Toeplitz operator also remembers the symbol form.  Here the two forms have much
more structure.  Outside a fixed ball their weights differ by
\[
 q(z)=\left(\frac{|z|}{1+|z|}\right)^\alpha,
 \qquad q(z)\longrightarrow1,
\]
while for \(\alpha>0\) all terms involving the low Taylor projection \(P\)
have finite rank.  Contributions from a fixed ball are better still: analytic
Taylor approximation makes them belong to every Schatten class.

The operator \(S\) changing the modified inner product into the canonical one
therefore satisfies \(S=I+K\) with \(K\) compact on every \(\Fc p\alpha\).
Radiality makes \(S\) diagonal with strictly positive diagonal entries, so it
is injective; Fredholm theory makes it invertible.  After multiplication by
\(S\), the canonical Toeplitz operator is the source Toeplitz operator with
symbol \(q\mu\), modulo the local and finite-rank corrections.  Since
\(q\asymp1\) off a fixed ball, all ball-mass criteria transfer unchanged.
To handle the quasi-Banach Berezin clause, the canonical weight is smoothed
only on a fixed ball.  Generalized-Fock kernel localization for the smoothed
weight transfers back through a compactly supported norm change and yields
exponential off-diagonal decay for the \emph{exact} canonical kernel.

\section{The transfer argument}

\begin{lemma}[Compactly supported forms are super-smoothing]\label{lem:local}
Suppose \(\sigma\) is a finite positive measure supported in a fixed ball.
The operator representing \(C_\sigma\), with respect to either of the two
equivalent Hilbert structures, belongs to \(\Sc_t\) for every \(t>0\).
The corresponding operator on \(\Fc p\alpha\) is compact for every
\(1\leq p<\infty\).  The same conclusion holds for the high-high modified
form containing \(|z|^{-\alpha}\), after applying \(Q\) when \(\alpha>0\).
\end{lemma}

\begin{proof}
Let \(\Pi_N\) be the Taylor projection onto total degree at most \(N\).  Its
rank is \(O(N^n)\).  On a fixed ball, the normalized monomial estimates for a
Gaussian Fock norm give
\[
 \sup_{|z|\leq R}|(I-\Pi_N)f(z)|
 \leq A_R\frac{B_R^N}{\sqrt{N!}}\,\|f\|_{\Fc2\alpha}.              \tag{4.1}
\]
For the high part, multiplication by \(|z|^{-\alpha/2}\) is harmless once
\(N>\alpha/2\), because every omitted term vanishes to higher order at zero.
Consequently the approximation numbers of the local embedding decay faster
than any power after a rank \(O(N^n)\) truncation.  The embedding, and hence
its positive square, lies in every Schatten class.  The same Taylor
truncations converge in operator norm on \(\Fc p\alpha\), proving compactness.
\end{proof}

Define \(S\) first on polynomials by
\[
 M(Sf,g)=C(f,g).                                                    \tag{4.2}
\]

\begin{lemma}[The norm-change operator]\label{lem:S}
For every \(1\leq p<\infty\), \(S\) extends to a bounded invertible operator
on \(\Fc p\alpha\), and
\[
 S-I\text{ is compact on }\Fc p\alpha.                             \tag{4.3}
\]
On \(\Fc2\alpha\), the compact difference belongs to every Schatten class
after removal of the exterior source-Toeplitz term described below.
The kernels satisfy
\[
 K^c_z=S^{-1}K^m_z.                                                 \tag{4.4}
\]
\end{lemma}

\begin{proof}
Take a cutoff at \(|z|=R\geq1\).  On the exterior high-degree block,
\(C-M\) is the modified Toeplitz form with scalar symbol \(q-1\), where
\(q=(|z|/(1+|z|))^\alpha\).  This symbol is bounded there and tends to zero.
The compactness part of Theorem B in \cite{CIJ}, applied to \(|q-1|dv\),
shows that this exterior difference is compact on every \(\Fc p\alpha\).
The interior difference is compact by Lemma \ref{lem:local}.  When
\(\alpha>0\), the remaining blocks contain \(P\) and hence have finite rank.
This proves \((4.3)\) and boundedness.

Both radial forms make the monomials orthogonal.  Thus \(S\) is diagonal on
homogeneous polynomials, and every diagonal coefficient is a ratio of two
strictly positive moments.  Hence \(S\) is injective on every \(\Fc p\alpha\).
Since \(S=I+\text{compact}\), it is Fredholm of index zero and therefore
invertible.  Finally, for every \(f\),
\[
 f(z)=M(f,K^m_z)=M(Sf,K^c_z)=M(f,SK^c_z),
\]
which gives \((4.4)\).
\end{proof}

\begin{lemma}[Toeplitz transfer modulo negligible blocks]\label{lem:transfer}
Assume \(\sup_z\mu(B(z,r))<\infty\).  Put
\[
 d\nu(z)=\mathbf1_{\{|z|>R\}}
 \left(\frac{|z|}{1+|z|}\right)^\alpha d\mu(z).
\]
Then, on every \(\Fc p\alpha\),
\[
 S\Tc_\mu=\Tm_\nu+E_\mu,                                         \tag{4.5}
\]
where \(E_\mu\) is compact.  On \(\Fc2\alpha\),
\(E_\mu\in\bigcap_{t>0}\Sc_t\).  In particular, for every \(t>0\),
\[
 \Tc_\mu\in\Sc_t\quad\Longleftrightarrow\quad\Tm_\nu\in\Sc_t.\tag{4.6}
\]
\end{lemma}

\begin{proof}
For \(\alpha\leq0\), the exterior canonical form is exactly the modified
form with symbol \(\nu\), because
\[
 |z|^{-\alpha}\left(\frac{|z|}{1+|z|}\right)^\alpha
 =(1+|z|)^{-\alpha}.
\]
For \(\alpha>0\), the same identity gives equality of the high-high blocks.
Every discrepancy involving \(P\) has finite rank.  The part supported in
\(B(0,R)\) is covered by Lemma \ref{lem:local}.  Equation \((4.5)\) follows
from the defining identity \(M(S\Tc_\mu f,g)=C_\mu(f,g)\).
The Schatten assertion follows because the error is in every \(\Sc_t\) and
left multiplication by the bounded invertible operator \(S\) preserves every
Schatten ideal, including the quasi-ideals with \(t<1\).
\end{proof}

\begin{proof}[Proof of Theorem \ref{thm:main}]
We first prove necessity of \((2.4)\) without any kernel estimate.  Put
\(E_a(z)=e^{z\cdot\overline a}\).  Gaussian localization gives
\[
 \|E_a\|_{\Fc p\alpha}\asymp
 e^{|a|^2/2}(1+|a|)^{-\alpha/p}.                                  \tag{4.7}
\]
Use the analogous conjugate-exponent normalization, with the unweighted
Fock \(F^\infty\) norm when \(p=1\).  If \(f_a,g_a\) are the two normalized
exponentials, then \((2.2)\) and \(|w-a|<r\) give
\[
 |C(\Tc_\mu f_a,g_a)|=C_\mu(f_a,g_a)\gtrsim\mu(B(a,r)).             \tag{4.8}
\]
The canonical pairing is bounded by Holder's inequality, so boundedness of
\(\Tc_\mu\) makes the right side uniformly bounded.

Conversely, \((2.4)\) implies that \(\nu\) in Lemma \ref{lem:transfer} has
uniformly bounded ball masses.  Theorem A of \cite{CIJ} makes \(\Tm_\nu\)
bounded on every \(\Fc p\alpha\).  Equation \((4.5)\) and invertibility of
\(S\) prove boundedness of \(\Tc_\mu\).

If \((2.5)\) holds, the ball masses of \(\nu\) vanish, so source Theorem B
and \((4.5)\) make \(\Tc_\mu\) compact.  Conversely, compactness first gives
boundedness and hence \((2.4)\).  Then \((4.5)\) makes \(\Tm_\nu\) compact;
source Theorem B implies that the ball masses of \(\nu\), and therefore those
of \(\mu\) outside a fixed ball, vanish.  This proves \((2.5)\).

Finally, source Theorem C says
\[
 \Tm_\nu\in\Sc_t\quad\Longleftrightarrow\quad
 \{\nu(B(a_j,r))\}\in\ell^t.
\]
Since \(q\asymp1\) on the exterior and only finitely many lattice cells meet
the omitted ball, this sequence lies in \(\ell^t\) exactly when
\(\{\mu(B(a_j,r))\}\) does.  Lemma \ref{lem:transfer} proves \((2.6)\).
The equivalence with \(\mu(B(\mathord\cdot,r))\in L^t\) is the standard
bounded-overlap lattice discretization (valid also for \(0<t<1\)).
\end{proof}

\section{Canonical Berezin transform}

Put
\[
 \phi_c(z)=\frac{|z|^2}{2}+\frac{\alpha}{2}\log(1+|z|),
 \qquad e^{-2\phi_c(z)}=e^{-|z|^2}(1+|z|)^{-\alpha}.
\]

\begin{lemma}[Exact canonical off-diagonal decay]\label{lem:canonical-decay}
There are constants \(C,\varepsilon>0\), depending only on \(n\) and
\(\alpha\), such that
\[
 |K^c_z(w)|e^{-\phi_c(z)-\phi_c(w)}
 \leq C e^{-\varepsilon|z-w|},                                   \tag{5.1}
\]
and
\[
 K^c_z(z)e^{-2\phi_c(z)}\asymp1.                                 \tag{5.2}
\]
Consequently,
\[
 |k^c_z(w)|^2e^{-2\phi_c(w)}\leq C e^{-2\varepsilon|z-w|}.       \tag{5.3}
\]
\end{lemma}

\begin{proof}
We first remove the harmless nonsmoothness of \(\phi_c\) at the origin without
changing it at infinity.  For \(s>0\), write \(\phi_c(z)=F_0(|z|^2)\).  If
\(r=\sqrt{s}\), the tangential and radial eigenvalues of the complex Hessian
are, up to the fixed normalization of \(dd^c\),
\[
 F_0'(s)=\frac12+\frac{\alpha}{4r(1+r)},\qquad
 F_0'(s)+sF_0''(s)=\frac12+\frac{\alpha}{8r(1+r)^2}.              \tag{5.4}
\]
They are bounded above and below by positive constants for \(r\) sufficiently
large.  A standard radial interpolation now gives
\(\psi\in C^2(\Cn)\), equal to \(\phi_c\) outside a ball, such that
\[
 c\,\omega_0\leq dd^c\psi\leq C\,\omega_0.                      \tag{5.5}
\]
For completeness, set \(g_0(s)=sF_0'(s)\) on a sufficiently large exterior
interval.  Extend it smoothly to \(s=0\), with \(g(s)=\lambda s\) near zero,
so that \(0<c\leq g'(s)\leq C\), while retaining \(g=g_0\) on the exterior.
This is possible because both \(g_0'(s)\) and \(g_0(s)/s\) have such bounds
there.  Integrating \(\Psi'(s)=g(s)/s\) and choosing the additive constant at
the joining radius gives \(\psi(z)=\Psi(|z|^2)\); its Hessian eigenvalues are
\(g(s)/s\) and \(g'(s)\), proving \((5.5)\).

Let \(H_\psi\) be the generalized Fock Hilbert space with weight
\(e^{-2\psi}\), and let \(L_z\) be its reproducing kernel.  The generalized
Fock kernel estimate \cite[Lemmas 2.1--2.2]{IVW} gives
\[
 |L_z(w)|e^{-\psi(z)-\psi(w)}\leq C e^{-\varepsilon|z-w|},
 \qquad L_z(z)e^{-2\psi(z)}\asymp1.                              \tag{5.6}
\]
Since \(\psi=\phi_c\) off a compact set, the two Hilbert norms differ by a
bounded compactly supported multiplier.  More precisely, with
\[
 h=e^{2(\psi-\phi_c)}-1,
\]
we have
\[
 C(f,g)=\ip{Af}{g}_{\psi},\qquad A=I+T_h^\psi,                   \tag{5.7}
\]
where \(T_h^\psi=P_\psi M_h\).  The operator \(A\) is positive and boundedly
invertible, while \(T_h^\psi\) is a compactly supported operator of the type
in Lemma \ref{lem:local}.  Reproducing gives
\[
 K^c_z=A^{-1}L_z=L_z+BL_z,
 \qquad B=A^{-1}-I=-A^{-1}T_h^\psi.                              \tag{5.8}
\]

Projection contractivity, compact support of \(h\), and \((5.6)\) imply
\[
 \|T_h^\psi L_z\|_\psi
 \leq \|hL_z\|_{L^2(e^{-2\psi})}
 \leq C e^{\psi(z)}e^{-\varepsilon|z|}.                         \tag{5.9}
\]
Point evaluation and the diagonal estimate in \((5.6)\) therefore give
\[
 |(BL_z)(w)|\leq C e^{\psi(z)+\psi(w)}e^{-\varepsilon|z|}.       \tag{5.10}
\]
The operator \(B\) is self-adjoint, so its kernel is Hermitian; interchanging
\(z,w\) gives the same bound with \(e^{-\varepsilon|w|}\).  Taking the better
bound and using
\(\max\{|z|,|w|\}\geq |z-w|/2\) shows that the correction in \((5.8)\)
has the exponential off-diagonal bound in \((5.1)\).  Combining this with
\((5.6)\), and using that \(\psi-\phi_c\) is bounded, proves \((5.1)\).

Equivalent Hilbert norms give comparable squared norms for every point
evaluation.  Thus \((5.6)\) implies \((5.2)\), and division by
\(K^c_z(z)\) gives \((5.3)\).
\end{proof}

We record why no explicit formula for \(K^c\) is needed for the bounded and
vanishing clauses.  The normalized modified kernels \(k^m_z\) converge weakly
to zero as \(|z|\to\infty\).  Lemma \ref{lem:S} therefore yields
\[
 \|(S^{-1}-I)k^m_z\|_{M}\longrightarrow0.                         \tag{5.11}
\]
The source near-diagonal lower estimate for \(K^m\), together with point
evaluation and \((4.4)\), implies that for some \(r_0,c>0\),
\[
 |K^c_z(w)|\geq c\sqrt{K^c_z(z)K^c_w(w)}
 \quad\text{when }|w-z|<r_0\text{ and }|z|\text{ is large}.        \tag{5.12}
\]
Indeed, the difference in \((5.11)\), evaluated at \(w\), is
\(o(\sqrt{K^m_w(w)})\), uniformly in \(w\), while the source kernel has a
fixed lower constant on the ball.  Equivalent Hilbert norms give
\[
 K^c_z(z)\asymp K^m_z(z)\asymp(1+|z|)^\alpha e^{|z|^2}.            \tag{5.13}
\]
It follows from \((5.12)\)--\((5.13)\) that
\[
 \widetilde\mu_c(z)\gtrsim\mu(B(z,r_0))                           \tag{5.14}
\]
for large \(|z|\).  Thus boundedness or vanishing of the canonical Berezin
transform forces the corresponding ball condition.  The converses follow
from boundedness or compactness of \(\Tc_\mu\); normalized canonical kernels
are weakly null at infinity.

For \(t\geq1\), positivity and Jensen's inequality give
\[
 \widetilde\mu_c(z)^t\leq C((\Tc_\mu)^t k^c_z,k^c_z).
\]
The canonical continuous-frame measure
\[
 K^c_z(z)e^{-|z|^2}(1+|z|)^{-\alpha}\,dv(z)
\]
is comparable with \(dv(z)\) by \((5.3)\).  Integration proves
\(\|\widetilde\mu_c\|_{L^t}^t\lesssim\|\Tc_\mu\|_{\Sc_t}^t\).
The reverse implication follows from \((5.14)\), lattice discretization, and
\((2.6)\).

It remains only to prove the forward implication when \(0<t<1\).  Suppose
\(\Tc_\mu\in\Sc_t\).  By \((2.6)\), if
\(m_j=\mu(B(a_j,r))\), then \(\{m_j\}\in\ell^t\).  Choose a measurable
partition \(\{D_j\}\) of \(\Cn\) with \(D_j\subset B(a_j,r)\).  Lemma
\ref{lem:canonical-decay} gives
\[
 \widetilde\mu_c(z)
 \leq C\int_{\Cn}e^{-2\varepsilon|z-w|}\,d\mu(w)
 \leq C_r\sum_j m_j e^{-\varepsilon|z-a_j|}.                    \tag{5.15}
\]
Since \(0<t<1\), subadditivity and translation invariance yield
\[
 \int_{\Cn}\widetilde\mu_c(z)^t\,dv(z)
 \leq C_{r,t}\sum_j m_j^t
 \int_{\Cn}e^{-t\varepsilon|z-a_j|}\,dv(z)
 \leq C'_{r,t}\sum_jm_j^t<\infty.                               \tag{5.16}
\]
This proves \((2.7)\) throughout the quasi-Banach range and completes all
claims.

\section{Verification, scope, and novelty}

\begin{itemize}
\item The proof uses only the source Theorems A--C, its modified-kernel
near-diagonal estimate, elementary Fredholm theory, Taylor approximation, and
the generalized-Fock kernel estimate of Isralowitz--Virtanen--Wolf.
\item The potentially singular factor at the origin is always placed after
the high Taylor projection \(Q\); Lemma \ref{lem:local} explicitly uses the
resulting order of vanishing.
\item All assertions involving arbitrary \(p\) concern boundedness and
compactness.  Schatten ideals concern the canonical Hilbert space
\(\Fc2\alpha\).
\item The exact canonical kernel is proved exponentially localized by a
compactly supported smoothing of the weight.  The correction kernel decays in
one variable by support localization and in the other by self-adjointness.
\item A bounded search on 27 August 2026 used the exact source id/title and
the canonical/natural-inner-product Fock--Sobolev Toeplitz keywords.  No
matching transfer result was located.  A 2023 paper still uses the modified
inner product and repeats the canonical-kernel difficulty; this is evidence
of continued relevance, not proof of novelty.
\end{itemize}

\paragraph{Human-review recommendation.}
Focused expert review should audit the local-form compactness in Lemma
\ref{lem:local}, the \(\alpha>0\) block identity \((4.5)\), and the radial
interpolation/self-adjoint two-variable decay in Lemma
\ref{lem:canonical-decay}.

\begin{thebibliography}{9}
\bibitem{CIJ}
H. R. Cho, J. Isralowitz and J.-C. Joo,
\emph{Toeplitz operators on Fock--Sobolev type spaces},
Integral Equations Operator Theory 82 (2015), 1--32;
\href{https://arxiv.org/abs/1307.4829}{arXiv:1307.4829}.

\bibitem{CCK}
H. R. Cho, B. R. Choe and H. Koo,
\emph{Fock--Sobolev spaces of fractional order},
Potential Analysis 43 (2015), 199--240;
\href{https://arxiv.org/abs/1307.2670}{arXiv:1307.2670}.

\bibitem{IVW}
J. Isralowitz, J. A. Virtanen and L. Wolf,
\emph{Schatten class Toeplitz operators on generalized Fock spaces},
J. Math. Anal. Appl. 421 (2015), 329--337;
\href{https://arxiv.org/abs/1402.2567}{arXiv:1402.2567}.
\end{thebibliography}

\end{document}
